\documentclass[11pt,a4paper]{article}
\usepackage[T1]{fontenc}
\usepackage[utf8]{inputenc}
\usepackage{lmodern}
\usepackage{microtype}
\usepackage[a4paper,margin=24mm,headheight=14pt]{geometry}
\usepackage{amsmath,amssymb,amsthm,mathtools,mathrsfs}
\usepackage{booktabs,tabularx,array}
\usepackage{enumitem}
\usepackage{fancyhdr}
\usepackage{needspace}
\usepackage[hidelinks,pdfauthor={},pdftitle={Exact Incremental SMILES Generation as an Online Constraint Problem}]{hyperref}
\usepackage[nameinlink,noabbrev]{cleveref}
\setlength{\parindent}{0pt}
\setlength{\parskip}{5.2pt plus 1pt minus 1pt}
\setlength{\emergencystretch}{2em}
\clubpenalty=10000
\widowpenalty=10000
\displaywidowpenalty=10000
\setlist{nosep,leftmargin=1.6em,topsep=4pt}
\renewcommand{\arraystretch}{1.12}
\newcolumntype{Y}{>{\raggedright\arraybackslash}X}
\pagestyle{fancy}
\fancyhf{}
\fancyhead[L]{\small Exact incremental SMILES}
\fancyhead[R]{\small Residual constraints and supported commitments}
\fancyfoot[C]{\thepage}
\renewcommand{\headrulewidth}{0.3pt}
\newtheorem{proposition}{Proposition}
\newtheorem{hypothesis}{Hypothesis}
\newcommand{\Sol}{\operatorname{Sol}}
\newcommand{\Supp}{\operatorname{Supp}}
\newcommand{\Accept}{\operatorname{Accept}}
\newcommand{\obs}{\operatorname{obs}}
\newcommand{\tw}{\operatorname{tw}}
\newcommand{\sgn}{\operatorname{sgn}}
\newcommand{\scope}{\operatorname{scope}}
\newcommand{\Inc}{\operatorname{Inc}}
\newcommand{\ind}[1]{\left[#1\right]}
\newcommand{\R}{\mathcal R}
\newcommand{\F}{\mathcal F}
\newcommand{\cC}{\mathcal C}
\newcommand{\cL}{\mathcal L}
\newcommand{\ttt}[1]{\texttt{#1}}
\newcommand{\none}{\bot}
\newcommand{\softO}{\widetilde O}
\begin{document}
\thispagestyle{plain}
\begin{center}
{\LARGE\bfseries Exact Incremental SMILES Generation\par}
\vspace{5pt}
{\large An online constraint problem over physical commitments\par}
\vspace{8pt}
{\small 24 October 2026\par}
\end{center}
\vspace{3pt}

A writer for a fixed molecule must offer exactly those next decisions that still
permit a correct completion. The objective is to retain the consequences of
unresolved choices as constraints, update those constraints after an actual
commitment, and enumerate supported decisions without constructing hypothetical
continuation runs. Constraint elimination may reason substantially about the
future; it must preserve the correlations needed by later commitments.

\begin{center}
\begin{tabularx}{\linewidth}{@{}lY@{}}
\toprule
\(x\) & The \emph{actual configuration}: what has been committed and the current control state.\\
\(z\) & One possible \emph{assignment} to unresolved constraint variables; not another writer state.\\
\(H_G(x)\) & The stored factors, summaries, and indexes; not a list of all assignments or completions.\\
\bottomrule
\end{tabularx}
\end{center}

\section{Fixed molecule, actual machine, exact support}
\subsection{The input and the target language}
Let \(G=(V,E,\partial,\alpha,\beta,\mathcal A,\mathcal S)\) be a finite molecular
graph. The sets \(V,E\) contain distinct physical atom and bond identifiers;
\(\partial(e)=\{u,v\}\) gives a bond's endpoints. Atom records \(\alpha\), bond
records \(\beta\), aromatic/delocalization records \(\mathcal A\), and stereo
records \(\mathcal S\) are fixed input. Physical identifiers are not SMILES atom
classes and need not appear in the emitted string.

The syntax and identity rules are fixed before choosing an algorithm. They
preserve graph atoms, connectivity, element, isotope, charge, numeric atom class,
hydrogen inventory, permitted bonding equivalences, and specified versus
unspecified stereo. Explicit hydrogen vertices remain distinct from non-graph
hydrogen counts; absent isotope is distinct from isotope zero. Aromatic/Kekul\'e
equivalence requires explicit rules, not whatever a toolkit happens to
normalize.\cite{note126,note127}

A \emph{labelled serialization} is a pair \((s,\lambda)\), where \(s\) is a finite
valid SMILES string for \(G\) and \(\lambda\) maps its atom occurrences bijectively
to \(V\), preserving the fixed facts. Write \(\cL_G\) for all such pairs. Equal
strings with different physical correspondences remain different labelled
outputs. The permitted unbounded leading-zero padding in isotope and class fields
can make \(\cL_G\) infinite even for finite \(G\); every accepting run is nevertheless
finite.\cite{note127}

\subsection{A local-commitment machine}
A machine supplies an initial state \(x_0\), a finite set \(A_G(x)\) of locally
admissible actions at each state, a deterministic update \(T_G(x,a)\), an
acceptance predicate \(\Accept_G(x)\), and an emission
\(\operatorname{emit}_G(x,a)\in\Sigma^*\). Here \(\Sigma\) is the character alphabet
and \(\Sigma^*\) includes the empty string. An action makes a primitive physical,
control, or lexical commitment; its permission rules do not call a completion
oracle. The same emitted text may accompany different actions and successors.

\Needspace{9\baselineskip}
An actual configuration has the schema
\[
  x=(C,\kappa,O,B,L).
\]
The record \(C\) contains committed atom/edge representations and registered
incidences; \(\kappa\) is the current owner and grammar phase; \(O\) is the actual
ring table with bound physical endpoints; \(B\) is the ordered stack of saved
branch-return records; and \(L\) is the current lexical residual, such as a
partially written atom or number. Exact record types and primitive actions are
part of a construction. Solver propagation and cache changes do not create new
actual commitments. This separates physical state from its maintained constraint
representation.\cite{note126}

A \emph{run} is a sequence of actions applied through \(T_G\); its output concatenates
the emissions and retains the physical atom correspondence. The machine is
\emph{adequate} if its finite accepting runs produce exactly \(\cL_G\). Primitive
actions cannot be whole future serializations or private solver steps disguised
as writer progress. A character-stream interface is a separate refinement: it
must preserve hidden physical distinctions and account for non-emitting steps.

Let \(\cC_G(x)\) be the set of finite accepting continuations from \(x\). It includes
the empty continuation exactly when \(\Accept_G(x)\) holds. The required next-choice
set is
\begin{equation}\label{eq:support}
 \boxed{\Supp_G(x)=\{a\in A_G(x):\cC_G(T_G(x,a))\ne\varnothing\}.}
\end{equation}
This is a specification, not an instruction to enumerate \(\cC_G(x)\).
Acceptance answers whether stopping \emph{now} is valid; support answers whether
taking an action still permits eventual acceptance. Existence of a finite
completion does not imply termination along every infinite sequence of choices.

\section{The residual CSP: relation, not a list of futures}
\subsection{Variables and factors}
For fixed \(G,x\), a finite-domain constraint satisfaction problem (CSP) is
\[
 F_x=(U_x,\{D_u\}_{u\in U_x},\F_x).
\]
A variable \(u\in U_x\) has a declared meaning and a finite domain \(D_u\). A factor
\(f\in\F_x\) has a scope \(S_f\subseteq U_x\) and a permitted relation
\(R_f\subseteq\prod_{u\in S_f}D_u\). An assignment \(z\) chooses one value
\(z(u)\in D_u\) for every variable. Its restriction to \(S_f\) is written
\(z|_{S_f}\). Define
\begin{equation}\label{eq:solutions}
 \Sol(F_x)=\{z\in\textstyle\prod_{u\in U_x}D_u:
                    z|_{S_f}\in R_f\text{ for every }f\in\F_x\}.
\end{equation}
The CSP is \emph{feasible} when this set is nonempty. Unary domain factors are
included; zero-variable true and false relations are allowed.

For relations \(P\) on variables \(A\) and \(Q\) on variables \(B\), define
\[
 P\Join Q=\{z\text{ on }A\cup B:z|_A\in P,\ z|_B\in Q\},
 \qquad \pi_Y(P)=\{z|_Y:z\in P\}\quad(Y\subseteq A).
\]
These are \emph{natural join} and \emph{projection}. Thus
\(\Sol(F_x)=\Join_{f\in\F_x}R_f\). Disjoint factors may remain separate: their
Cartesian product need not be expanded. Shared scopes require agreement and
preservation of correlations. For instance, replacing \(u\ne v\) on
\(\{0,1\}^2\) by the two marginal domains incorrectly admits \((0,0)\) and
\((1,1)\). Disconnected molecular components are not automatically independent
when they share control or ring resources.\cite{note127}

\subsection{What is represented exactly?}
Let \(W_G\) be the fixed physical observation variables: normalized atom lexemes,
local incidence-registration orders, and per-edge direct/ring representation
records. A lexeme is a token spelling; normalization removes only the designated
unbounded padding, not other permitted spelling alternatives. A complete labelled
serialization induces an assignment \(\obs_G(s,\lambda)\) to \(W_G\). Different
global orders and padding lengths may have the same observations.\cite{note127}

Let \(\R_G(x)\) contain the observation tuples of all complete labelled
serializations extending the actual commitments and control in \(x\).
An \emph{exact physical factorization} has variables \(U_x=W_G\cup Y_x\), where
\(Y_x\) are explicitly defined auxiliaries, and satisfies
\begin{equation}\label{eq:exact}
 \boxed{\pi_{W_G}\Sol(F_x)=\R_G(x).}
\end{equation}
Consequently \(\Sol(F_x)\ne\varnothing\) exactly when \(x\) is completable.
Equation~\eqref{eq:exact} preserves the joint physical-choice relation, not merely
its nonemptiness. It is stronger than the behavioral requirement
\eqref{eq:support}; failure of this particular representation requirement would
not rule out every exact generator.\cite{note127}

\section{Concrete molecular variables and factors}
\subsection{Atom and edge representations}
An atom record fixes
\[
 \alpha(v)=(Z_v,\iota_v,q_v,h_v,c_v),
\]
where \(Z_v\) is the element (or a permitted wildcard), \(\iota_v\) is an optional
isotope value, \(q_v\) is charge, \(h_v\) is the non-graph hydrogen count, and
\(c_v\) is the numeric atom class. Let \(\Inc(v)\) be its incident physical edges.
The following variables choose an encoding of these fixed facts, not a different
molecule:

\begin{center}
\begin{tabularx}{\linewidth}{@{}lY@{}}
\toprule
Variable & Meaning\\
\midrule
\(L_v\) & Normalized atom token, including bracket choice, numeric fields, hydrogen spelling, and stereo mark.\\
\(\Pi_v\) & A permutation of \(\Inc(v)\): its incidence-registration order.\\
\(R_e\) & A direct/ring record. Direct: introducing endpoint, branch/inline attachment, bond symbol. Ring: opening endpoint, label, endpoint symbols and numeral spellings.\\
\(b_e\) & Auxiliary decoded bond type for edge \(e\).\\
\(d_{v,e}\) & Auxiliary directional mark viewed outward from \(v\): \(0,+1,-1\), with zero meaning unmarked.\\
\(O_t\) & Ordered stereo carriers for stereo record \(t\), including any permitted non-graph placeholders.\\
\bottomrule
\end{tabularx}
\end{center}

An atom factor
\(A_v(L_v,(b_e)_{e\in\Inc(v)})\) accepts exactly the token/bond combinations
that decode to \(\alpha(v)\). Explicit bracket hydrogen counts and shorthand
implicit hydrogens are different decoding cases; no additional generic
``chemically plausible valence'' filter is imposed. An edge factor
\(E_e(L_u,L_v,R_e,b_e,d_{u,e},d_{v,e})\), for \(\partial(e)=\{u,v\}\), decodes
its endpoint symbols consistently. For fixed nonaromatic bonding add
\(\ind{b_e=\beta(e)}\), where \(\ind{P}\) is one if proposition \(P\) is true and
zero otherwise. Aromatic alternatives require joint factors rather than this
single equality.\cite{opensmiles,note127}

\subsection{Stereo is a constraint on named carriers and their order}
A stereo record is
\((\mathrm{kind},\mathrm{focus},P_t,c_t^\star)\): the geometry, its atom/bond/axis
focus, a reference list \(P_t\) of identified carriers, and its fixed target
configuration \(c_t^\star\). An ordinary carrier is a physical bond incidence;
a counted hydrogen or other prescribed implicit carrier is a distinct
placeholder, not an invented graph vertex. The value \(\none\) denotes
\emph{unspecified} configuration. It is not a variable ranging over the specified
configurations.

For a specified tetrahedral centre \(v\), let \(P_v=(p_1,p_2,p_3,p_4)\) and let
\(\sigma_v\in\{+1,-1\}\) be the target orientation relative to this reference.
Let \(\pi(O_v)\) be the permutation taking \(P_v\) to the carrier order \(O_v\).
Choose the orientation convention \(\chi(\ttt{@})=+1\),
\(\chi(\ttt{@@})=-1\), and normalize equivalent long-form marks when evaluating
\(\chi\). With \(m(L_v)\) the mark decoded from the token, the factor is
\begin{equation}\label{eq:tetra}
 T_v(O_v,L_v)=
 \ind{\chi(m(L_v))\sgn(\pi(O_v))=\sigma_v}.
\end{equation}
Here \(\sgn\) is permutation parity; an exchange of two carriers reverses its
sign. The target configuration belongs to the molecule; \ttt{@} and \ttt{@@}
belong to its order-dependent encoding.\cite{opensmiles}

For reference carriers \((F,Cl,Br,I)\) and \(\sigma_v=+1\),
\ttt{F[C@](Cl)(Br)I} and \ttt{F[C@@](Br)(Cl)I} obey the same factor.
The second exchanges two carriers and also reverses the mark. A registration
factor links \(O_v\) to \(\Pi_v\), the edge records, and the fixed rules for
incoming attachments, implicit carriers, and ring references. The same order
variables therefore participate in chemistry and control; testing their
feasibility separately is insufficient.

For an ordinary stereogenic double bond \(e=uv\), fix two non-axis carriers at
each endpoint, numbered \(0,1\). Let \(d_0,d_1\in\{0,+1,-1\}\) be their written
bond directions viewed outward from that endpoint: \(+1\) for a slash and
\(-1\) for a backslash. Reversing the view reverses the sign; an unwritten
hydrogen carrier has value zero. Introduce a reference-side value \(r\in\{\none,+1,-1\}\) with
\[
 \operatorname{Side}(d_0,d_1,r)\iff
 \begin{cases}
   d_0=d_1=0, &r=\none,\\
   (d_0,d_1)\in\{(r,0),(0,-r),(r,-r)\}, &r\in\{+1,-1\}.
 \end{cases}
\]
For a specified target \(\tau_e\), with \(+1\) meaning the reference carriers are
on the same side and \(-1\) opposite, require
\[
 r_u\ne\none,\qquad r_v\ne\none,\qquad r_ur_v=\tau_e.
\]
A physical directional bond supplies shared variables with
\(d_{u,e}=-d_{v,e}\); adjacent stereo systems cannot choose its direction
independently. For an unspecified target, the decoded result must remain
\(\none\), not become an existentially chosen specified result. This is a
finite-table encoding of the ordinary two-ended convention, not a treatment of
all extended stereo classes.

Other geometries use a finite relation
\[
 T_t(O_t,L_{\mathrm{focus}(t)})=
 \ind{\operatorname{Decode}_{\mathrm{kind}(t)}(O_t,m(L_{\mathrm{focus}(t)}))
             =c_t^\star}.
\]
For an axis-centred geometry the token variable is the one designated by its
syntax. Each decoder must be an explicit geometry/permutation table, including
its unspecified value; parity alone does not describe every geometry.
Extended systems may involve terminal carriers beyond the immediate neighbours
of the marked atom. Their factor scopes must be counted, not assumed local
because the stereo mark is local.\cite{note126}

\subsection{Aromatic and structural coupling}
For a benzene \emph{Kekul\'e assignment}, let \(k_i\in\{0,1\}\) encode a single
or double cyclic edge. The six binary factors
\[
 \ind{k_{i-1}+k_i=1}\qquad(i\bmod 6)
\]
permit exactly the two alternating assignments. Independent single/double
edge domains would admit invalid combinations. General aromatic systems need
the declared bonding-and-valence relations for their element, charge, hydrogen,
and external-bond context. These equations are an explicit example, not a
complete aromaticity specification.

\Needspace{8\baselineskip}
The complete constraint system has the form
\begin{equation}\label{eq:blocks}
 \Phi_{G,x}=\Phi_G^{\mathrm{chem}}\land
             \Phi_x^{\mathrm{commit}}\land
             \Phi_{G,x}^{\mathrm{struct}}.
\end{equation}
The blocks are conjunctions of scoped factors. Chemical factors preserve the
fixed facts; commitment factors restrict choices already made; structural
factors must express joint realizability under introductions, branch control,
ring reuse, lexical progress, and incidence registration. A necessary local
factor, for example, permits at most one incident edge to introduce a given
atom directly. It does not establish global realizability. Ring labels are
not globally all-different: only simultaneous occupancy is exclusive.
Constructing the exact structural factors is part of the question, not a supplied
predicate named ``a serialization exists.''\cite{note127}

\section{Online reduction and safe forgetting}
\subsection{What changes when a choice is made?}
At selected step \(t\), retain \(x_t,F_t,H_t\). Let \(\delta_t\) record only
changed actual fields. An action \(a_t\) gives
\[
 x_{t+1}=T_G(x_t,a_t),\qquad
 F_{t+1}=\operatorname{Edit}(F_t,\varepsilon_t(a_t)),
\]
where \(\varepsilon_t(a_t)\) is the generated collection of base-factor edits.
It may create/delete variables, change domains, or insert/delete/replace
factors. A mark commitment restricts \(L_v\); an incidence event restricts
\(\Pi_v\); a ring opening fixes parts of \(R_e\) and changes actual slot records;
a phase transition replaces applicable control constraints. The target molecule
is known from the start: unvisited atoms' constraints are not postponed until
arrival. An actual commitment restricts compatible complete runs; replacing
factors does not authorize new completed encodings.\cite{note127}

\emph{Conditioning} incorporates a decision, for example by adding \(u=d\).
\emph{Propagation} derives consequences without making another decision.
\emph{Elimination} removes variables while retaining their joint effect:
\begin{equation}\label{eq:elim}
 M(B)=\exists I\;\bigwedge_{f\in\F_{\rm region}} f.
\end{equation}
Here \(I\) are internal variables and \(B\) are retained boundary variables;
factors in \eqref{eq:elim} are identified with their Boolean membership
predicates. The relation \(M\) is a join followed by projection. It can be much
larger than the original factors. A fixed point of incomplete domain propagation
is not a certificate of global feasibility.\cite{poole}

For a tetrahedral centre with its incoming carrier fixed, let the remaining
carriers be \(a,b,c\), and choose the reference sign so \ttt{@} requires an even
permutation. Committing \ttt{@} leaves
\(O\in\{(a,b,c),(b,c,a),(c,a,b)\}\). Registering \(b\) next leaves only
\(O=(b,c,a)\). The remaining stereo consequence is ``register \(c\), then \(a\).''
This deduction does not perform either action, and the structural factors must
still permit them. The fixed target stereo has not changed.

A feasibility summary can answer the current query correctly yet lose
information needed by the next commitment. For example, \(\exists m\,T_v(O,m)\) permits every carrier permutation, because each
has a compatible mark. Committing \ttt{@} later cannot recover the required
parity from that projection alone. Either keep \(m\) exposed, or retain the
original factor and its dependencies so the affected region can be reopened,
restricted, and summarized again.

\begin{proposition}[Safe forgetting]\label{prop:forget}
Let \(R(I,B)\) be a relation whose internal variables \(I\) do not occur in a
later constraint \(C(B)\). Then
\begin{equation}\label{eq:forget}
 \exists I\,[R(I,B)\land C(B)]
   =\bigl(\exists I\,R(I,B)\bigr)\land C(B).
\end{equation}
Consequently its boundary summary is sufficient for that later constraint.
If a later edit touches \(I\), this guarantee does not apply.
\end{proposition}
\begin{proof}
For a fixed assignment to \(B\), the truth of \(C\) is independent of the
existentially quantified values of \(I\).
\end{proof}

``Already written'' therefore does not mean ``safe to discard.'' An atom can
remain relevant through an open ring, a saved return, or a stereo correlation.
A usable summary exposes every dependency through which outside constraints
interact with its interior; otherwise an edit requires reopening that interior.
Actual grammar state and base factors must remain available: the next projected
relation need not be obtainable by conditioning the previous projection alone.
Under \eqref{eq:exact}, eliminating physical observation variables from a query
cache must retain their joint reconstruction factors or their committed values;
otherwise the full physical relation has been lost. No single finite summary
need encode every future interaction sequence in advance.\cite{note126,note127}

\subsection{The online support query}
For any immediate candidate action, form its successor system semantically as
\(F_t^{\,a}=\operatorname{Edit}(F_t,\varepsilon_t(a))\). Exactness gives
\begin{equation}\label{eq:online}
 \boxed{\Supp_G(x_t)=
   \{a\in A_G(x_t):\Sol(F_t^{\,a})\ne\varnothing\}.}
\end{equation}
When an action only sets \(u=d\), support is membership of \(d\) in the projection
of the \emph{whole} solution relation onto \(u\). General actions replace factors
or control records. Equation~\eqref{eq:online} does not require a separate full
solve for every candidate.

The interface is \(\operatorname{Init}(G)\), \(\operatorname{List}(H_G(x))\), and
\(\operatorname{Select}(a)\). Initialization detects an empty target language.
Listing reports acceptance and enumerates every supported primitive action
once, with a successor handle or compact delta. Selection installs that
already-feasible successor without repeating a complete support decision.
Tentative candidates are reversible or persistent; their rejection costs belong
to listing. Multiple satisfying assignments are witnesses, not duplicate machine
outputs.\cite{note127}

\section{A conditional support-enumeration kernel}
Suppose a residual CSP has a fixed balanced binary composition tree with \(n\)
leaf modules, height \(O(\log n)\), domain size at most \(d\), and at most \(w\)
boundary variables per region. A leaf module is a bounded local group of factors.
Assume every immediate action has one unique owning leaf and replaces only that
module, with a bounded number of alternatives per leaf. Current factors and
\emph{all candidate alternatives} must fit the same variable/interface envelope.
An actual update changes \(r\) leaves, including changes in candidate availability.
These are explicit sufficient conditions, not established SMILES properties.

At a region \(t\) with boundary \(B_t\), retain \(M_t^0\), its feasible boundary
assignments without an edit, and \(M_t^1\), its feasible boundary assignments
after exactly one candidate edit somewhere inside. Leaf relations are computed
from the current module and the union of its candidate replacements. For
children \(l,r\), composition is
\begin{align}
 M_t^0&=\pi_{B_t}(M_l^0\Join M_r^0),\label{eq:zero}\\
 M_t^1&=\pi_{B_t}\bigl((M_l^1\Join M_r^0)
                    \cup(M_l^0\Join M_r^1)\bigr).\label{eq:one}
\end{align}
The single edit is in the left or right child; the other child stays unchanged.
This summarizes immediate actions, not sequences of hypothetical writer steps.

During enumeration, carry the outside-feasibility relation \(E_t\) on \(B_t\),
with the outside unedited. Skip region \(t\) when
\(M_t^1\Join E_t=\varnothing\). When descending left, use
\[
 E_l=\pi_{B_l}(E_t\Join M_r^0),
\]
and symmetrically for the right child. At a leaf, test its bounded alternative
set and emit supported actions. The traversal follows action owners, not
assignments, so different witnesses cannot duplicate an output. Outside
relations are computed along queried paths rather than refreshed globally.

Explicit boundary tables yield preprocessing/storage \(O(nd^{O(w)})\), maintenance
\(O(rd^{O(w)}\log n)\), and initial, per-answer, and terminal delay
\(O(d^{O(w)}\log n)\), plus output size. Between outputs only logarithmically
many tree nodes need be crossed: empty subtrees are rejected at their roots.
No claim follows when actions affect distant modules, candidate availability
changes globally, or the common envelope has large width. Small factor arity
alone does not supply the missing interface bound.

\section{A structural fragment and its obstruction}\label{sec:fragment}
\subsection{Nonbinding ring capacity}
Consider the fixed-chemistry structural model of the working derivation, with
no stereo or previously imposed incidence-order restrictions. The graph is
simple; actions complete atom tokens or ring endpoints. A first ring endpoint
binds its physical edge. Dot introductions, dot-prefixed branches, inline
continuations, arbitrary nesting, and ring-label reuse are retained.
The grammar places an atom's ring references before its branches; returning
from a branch does not reopen its ring phase.\cite{opensmiles,attempt}

Each atom is unentered (\(\mathsf U\)), current in ring phase (\(\mathsf R\)),
current or saved in branch phase (\(\mathsf B\)), or permanently departed
(\(\mathsf D\)). Edges are unrepresented, represented directly, open rings, or
closed rings. For an unentered atom \(v\), define
\[
 f_x(v)=\#\{u:\ u\text{ has status }\mathsf B,
                      \ uv\text{ is unrepresented}\}.
\]
An unrepresented edge from a \(\mathsf B\) owner must introduce its other endpoint
directly: the owner can no longer emit a ring reference. Thus \(f_x(v)>1\)
requires two different introductions of the same physical atom.

\begin{proposition}[Exact local criterion for the structural fragment]
Assume \(K\ge |E|\), so ring capacity cannot bind. A locally reachable
configuration is completable exactly when \(f_x(v)\le1\) for every unentered
atom and all edges satisfy the following availability test: an unrepresented
edge has endpoint statuses \(\mathsf U\!\! -\!\mathsf U\),
\(\mathsf R\!\! -\!\mathsf U\), or \(\mathsf B\!\! -\!\mathsf U\), in either
order; an open ring's remaining endpoint is unentered or current in ring phase;
and completed representations are valid.
\end{proposition}
\begin{proof}
Necessity follows from unique introduction and the loss of endpoint opportunities.
For sufficiency, enter an initial atom if necessary. At each current ring-phase atom,
close incoming rings and open
all its still-unrepresented edges to unentered atoms; capacity suffices. At each
branch-phase owner, introduce its forced children as bonded branches. Each child
performs the ring procedure and can return immediately, since no other owner
needs to introduce it. Discharge saved owners in stack order. At the outermost
owner, introduce remaining atoms through dots. Every ring eventually closes and
all atoms are consumed. This proves existence; the runtime need not execute
this particular completion strategy.\cite{attempt}
\end{proof}

Maintain edge violations, forced-parent counts, unrepresented-incidence counts,
and the indexed set
\(\operatorname{Free}(x)=\{v:\ v\text{ is unentered},\ f_x(v)=0\}\).
At maximum degree \(\Delta\), actual changes affect \(O(\Delta)\) incidences.
In branch phase, supported dot branches choose directly from
\(\operatorname{Free}(x)\); an inline dot additionally requires no remaining
unrepresented incidence at the departing owner. Every unrepresented incident
edge supports bonded branch entry; bonded inline entry requires that it be the
owner's last such edge. This gives local maintenance and indexed enumeration,
not a scan of trial completions.\cite{attempt}

\subsection{Binding capacity requires joint resource--control information}
Take the complete graph on \(A,B,C,D\), with \(K=2\). Let \(A\) be a saved
branch-phase owner, \(B\) current in ring phase, and \(C,D\) unentered. The edge
\(AB\) is direct; \(AC\) and \(BC\) occupy both ring slots; \(AD,BD,CD\) are
unrepresented. The local criterion passes. Nevertheless, \(AD\) forces \(D\)
to be introduced from \(A\), so \(BD\) must open as a ring before \(B\) leaves
ring phase. Both slots are released only at \(C\), but entering \(C\) irreversibly
ends that phase. The required opportunity has expired before the resource is
available.

The parent with only \(AC\) open is completable: opening \(BD\) instead of
\(BC\), entering \(C\) directly from \(B\), closing \(AC\), opening \(CD\),
returning to \(A\), and introducing \(D\) completes the graph within two slots.
Thus the obstruction is a dead candidate from a live parent, not merely an
initially impossible graph. Occupancy counts and forced-parent counts alone
cannot distinguish the alternatives. Exact factors must retain which endpoint
releases a slot, which owner must introduce an atom, and which phase change
removes an endpoint opportunity.\cite{attempt}

\section{Representation restrictions and the research hypothesis}
\subsection{An explicit, falsifiable cost contract}
Fix parameters \(\theta=(\mathscr D,\Delta,K,b)\): a declared chemical/lexical
relation vocabulary \(\mathscr D\), maximum physical degree \(\Delta\), ring
capacity \(K\), and maximum bit length \(b\) of input numeric values. Standard
numeric ring labels give \(K=100\); other budgets define analysis variants, not
new standard syntax. The parameter \(b\) does not bound output padding. A uniform finite library supplies
bounded-arity local factor schemas, with explicit parameter-bounded evaluation
costs. Variables are owned by physical atoms/edges or actual ring slots, return
frames, and lexical records. Per-owner slot counts and local domain sizes are
parameter-bounded; physical identifiers and large numeric values still cost
space. Each primitive actual change \(\delta\) generates at most
\(O_\theta(1+|\delta|)\) base edits, with their construction charged.\cite{note127}

No primitive factor may call a completion oracle. Storage may not be indexed
by full future executions, future visited-atom subsets, or global future visit
and branch-end schedules. Larger derived relations are permitted, with all
indices, dependencies, construction, and recomputation charged. Ordinary
physical bookkeeping may remain procedural. Already-built independent components
must combine additively without expanding their Cartesian product. These
restrictions define a proposed architecture, not a claim that all exact
generators must use it.\cite{note126,note127}

Let \(N\) be encoded input size, \(\ell\) the number of selected primitive actions,
\(P\) the number of actual emitted bytes, \(Q\) the number of listing calls,
\(R\) the total size of returned successor records, and
\(A=\sum_t|\delta_t|\) the total encoded size of changed actual records.
Let \(M\) be current actual-state storage and \(R_{\rm live}\) retained successor
storage. Set \(\Lambda=\log(2+N+\ell+P+Q+R+A)\). Use a word RAM with word size
proportional to \(\Lambda\); multiword values require multiple operations.
All displayed work and storage bounds use this fixed accounting convention.

The strong one-walk target asks for a computable function \(f\) and a constant
\(c\), independent of the input and run, such that
\begin{align}
 \text{preprocessing}&\le f(\theta)N\log^c(2+N),\nonumber\\
 \text{total work}&\le f(\theta)(N+Q+R+\ell+P+A)\Lambda^c,\label{eq:cost}\\
 \text{live storage}&\le f(\theta)(N+M+R_{\rm live})\Lambda^c.\nonumber
\end{align}
After a selected edit, maintenance is at most
\(f(\theta)(1+|\delta_t|)\Lambda^c\); each iterator call, including its first and
exhaustion, costs at most \(f(\theta)(1+B)\Lambda^c\), where \(B\) is the returned
record/payload size and is zero for exhaustion. Rejected trials and derived
propagation are work, not actual changes or outputs. An optional transcript is
charged separately in storage. Retaining a tree of successors requires an
additional persistence bound. These are targets, not established properties.

An additional independently defined interaction-width parameter gives a
qualified target \(f(\theta,w)\), not a proof of the width-independent one.
The width must describe the joint physical/control/factor structure and the
candidate-edit effects; it cannot be defined as measured running time. A
connectivity-width bound alone does not bound every possible factor scope.
\Cref{app:width} gives a relevant structural bound, but not the missing exact
constraint compiler.

\begin{hypothesis}[Incremental exact physical support]\label{hyp:main}
There exist an adequate local-commitment machine, a uniform physical-factor
compiler and update rules, and algorithms \(\operatorname{Init/List/Select}\)
such that, for every finite input graph and every finite walk selected through
supported actions, \eqref{eq:exact} holds initially and for every immediate
candidate successor, listing returns exactly \eqref{eq:support} and the correct
acceptance answer, and the representation and cost requirements above hold.
Feasibility reasoning beyond actual and immediate trial transitions operates
on physical constraint relations, not recursively advanced hypothetical writer
runs.
\end{hypothesis}

The unknown is both the exact encoding and its efficient maintenance. Small
factors do not ensure small elimination results, and fast feasibility tests for
individual candidates do not ensure small enumeration delay. A negative result
must identify whether it concerns this representation class, its cost targets,
or the broader behavioral problem. The central question is:
\begin{quote}
\textbf{Can the remaining physical choices be reduced exactly under actual
commitments, preserving the correlations needed by later choices, so that all
and only completable next decisions can be enumerated without expanding future
executions?}
\end{quote}

\appendix
\section{A bounded-width host for the structural model}\label{app:width}
A \emph{tree decomposition} is a tree of vertex bags such that every graph vertex
occurs, every graph edge has both endpoints in a bag, and the bags containing
any fixed vertex form a connected subtree. Its width is the largest bag size
minus one; \(\tw\) is the minimum such width. The following bounds concern the
simple-graph structural model in \cref{sec:fragment},
not arbitrary chemical factor scopes. They are the constructive statements
from the working derivation.\cite{attempt}

\begin{proposition}[Ring-budget host bounds]
If a complete structural serialization uses at most \(K\) simultaneously open
ring labels, then \(\tw(G)\le4K+5\). If \(x\) is completable and \(J_x\) augments
\(G\) by actual introduction links (including virtual dot links) and the path
through saved return owners ending at the current atom, then
\(\tw(J_x)\le4K+9\).
\end{proposition}
\begin{proof}
Connect every noninitial atom to the owner from which it is introduced, using
a virtual link for a dot. These links form an ordered introduction tree \(T\)
whose preorder is the atom-occurrence order. The other physical edges are ring
edges. At most \(K\) ring edges cross any cut in this order, since their first
endpoints have been emitted and their second endpoints have not.

Replace each high-child-count tree owner by a chain of copies, retaining its
ring incidences at the first copy. Attach successive children along the chain,
visiting each child before continuing along it. The expanded tree has at most
two children per node and preserves the original atom order. Contracting the
copy chains recovers \(T\).

Build a binary partition hierarchy by combining a node with its first complete
child subtree and then its second. Each region is a preorder interval with at
most three crossing tree edges. At most \(2K\) ring edges cross its two interval
ends. Thus each region has at most \(s=2K+3\) internal boundary vertices, where
a boundary vertex is one with a neighbour outside the region.

Use singleton bags at leaves and, at an internal hierarchy node, the union of
its two children's internal boundary sets. Bag size is at most \(2s\). Every
edge is covered at the first join of its endpoints' regions. A vertex remains
in consecutive ancestor bags until its region contains all its neighbours,
which proves connected occurrence. Contraction and deletion preserve a
no-larger decomposition, giving \(\tw(G)\le2s-1=4K+5\).

For \(J_x\), choose an accepting extension only to prove existence of the layout.
The actual introduction links belong to its introduction tree. Saved owners
and the current atom occur in increasing preorder, so their path adds at most
one crossing edge at any cut. Replace \(K\) by \(K+1\) in the construction.
\end{proof}

The proof does not instruct an implementation to find a completion before
constructing a decomposition. It establishes existence of a bounded-width host.
It neither supplies a bounded-width exact factor encoding on that host nor
bounds the cost of maintaining one under all candidate edits. Deep actual
stacks are compatible with the bound; unresolved stereo, resource, or control
relations must still have their scopes and payloads accounted for.

\begin{thebibliography}{9}
\small
\bibitem{note126}
\emph{Exact online SMILES generation without repeated future-execution search}.
Research question, 16 October 2026. Supplied file
\texttt{126\_exact\_incremental\_smiles\_feasibility\_research\_question.txt}.
Sections 1--5: molecular identity, actual-state semantics, joint feasibility,
and accounting requirements.

\bibitem{note127}
\emph{Factored online SMILES generation: an attempt at a formal problem}.
16 October 2026. Supplied file
\texttt{127\_factored\_online\_smiles\_formal\_problem.txt}.
Sections 3--8: relational semantics, exact projected invariant, physical-owner
restrictions, online edits, and proposed work bounds.

\bibitem{attempt}
\emph{Derivations: a local kernel, its boundary, and a width bound}.
Working derivation, supplied as \texttt{smiles\_kernel\_attempt/derivations.md}.
Parts A--G: structural model, nonbinding-capacity criterion, two-slot
counterexample, and physical/control host-width constructions.

\bibitem{opensmiles}
Craig A. James and contributors.
\emph{OpenSMILES specification}, version 1.0, 15 May 2016.
Formal Grammar; Reading SMILES: Atoms, Rings, Chirality, and Aromaticity.
\url{https://opensmiles.org/opensmiles.html}.
Source: \url{https://raw.githubusercontent.com/opensmiles/OpenSMILES/master/opensmiles.asciidoc}.

\bibitem{poole}
David L. Poole and Alan K. Mackworth.
\emph{Artificial Intelligence: Foundations of Computational Agents},
3rd ed., 2023. Section 4.5, Variable Elimination.
\url{https://artint.info/3e/html/ArtInt3e.Ch4.S5.html}.
\end{thebibliography}
\end{document}
